\section{Phase 0 Full Results}
\label{app:phase0}

Phase~0 reproduced the LeWM baselines on both environments and ran offset
sweeps, three-cost attribution, per-pair categorization, and aggregate latent
diagnostics.  These results motivated the audit framing developed in
Sections~3--4 of the main text.

\subsection{PushT Baseline and Offset Sweep}

The PushT baseline was reproduced at 98\% success on 50 episodes (paper
reference: 96\%; the difference is within finite-sample variation at $n=50$).
Table~\ref{tab:pusht-offset} reports the offset sweep.  Planning success drops
sharply beyond offset~25, falling from 96\% to 10\% at offset~100.  All runs
used seed~42 on Apple Silicon MPS; the full sweep took approximately one hour.
These numbers are plotted in Figure~\ref{fig:motivation}a.

\begin{table}[h]
\centering
\caption{PushT planning success versus goal offset.  Each row evaluates 50
episodes at the indicated offset (number of environment steps between the
initial observation and the goal observation in the expert demonstration).}
\label{tab:pusht-offset}
\begin{tabular}{lcccc}
\toprule
Goal offset (steps) & 25 & 50 & 75 & 100 \\
\midrule
Success rate (\%) & 96 & 58 & 16 & 10 \\
Successes / episodes & 48/50 & 29/50 & 8/50 & 5/50 \\
\bottomrule
\end{tabular}
\end{table}

\subsection{Cube Baseline and Offset Sweep}

The OGBench-Cube baseline was reproduced at 66\% on 50 episodes (paper
reference: 74\%).  Unlike PushT, Cube success does not degrade with goal
offset: it remains between 50\% and 68\% across all four offsets
(Table~\ref{tab:cube-offset}).  This flat curve indicates a
horizon-independent visual encoding difficulty rather than the
displacement-dependent breakdown observed in PushT.  Cube entered the audit
at Phase~2 with its own 100-pair stratified artifact (Appendix~\ref{app:stage1b}).

\begin{table}[h]
\centering
\caption{OGBench-Cube planning success versus goal offset.  Each row evaluates
50 episodes.  Success is defined by
$\|x_\text{cube}^\text{terminal} - x_\text{cube}^\text{goal}\|_2 \le 0.04$.}
\label{tab:cube-offset}
\begin{tabular}{lcccc}
\toprule
Goal offset (steps) & 25 & 50 & 75 & 100 \\
\midrule
Success rate (\%) & 68 & 50 & 58 & 50 \\
Successes / episodes & 34/50 & 25/50 & 29/50 & 25/50 \\
\bottomrule
\end{tabular}
\end{table}

\subsection{Three-Cost Attribution}
\label{app:three-cost}

At offset~50, 30 PushT diagnostic pairs were each evaluated with 40 candidate
actions (10 per source: expert data, random, CEM-early, CEM-late; 1{,}200
records total).  Three pairwise Spearman correlations decompose the end-to-end
cost alignment into a predictor-fidelity term and an encoder-physics term:

\begin{itemize}
\item \textbf{Predictor fidelity} $\rho(\Cmodel,\, C_\text{real\_z})$: how
  well the imagined terminal latent from the predictor tracks the real
  terminal latent.
\item \textbf{Encoder-physics alignment}
  $\rho(C_\text{real\_z},\, C_\text{real\_state})$: how well latent distance
  from the \emph{real} terminal encoding tracks physical task cost.
\item \textbf{End-to-end alignment}
  $\rho(\Cmodel,\, C_\text{real\_state})$: the product-path correlation that
  CEM planning ultimately depends on.
\end{itemize}

\noindent
Here $\Cmodel = \|\zhatH - \zgoal\|^2$ is the model cost used by CEM,
$C_\text{real\_z} = \|z_\text{terminal} - \zgoal\|^2$ uses the encoder
applied to the actual terminal observation, and
$C_\text{real\_state}$ is the V1 hinge cost combining block position error
and angle error.

Table~\ref{tab:three-cost} reports both global (all 1{,}200 records pooled)
and per-pair (mean $\pm$ SEM over 30 pairs) correlations.  Predictor fidelity
is high in both views.  Encoder-physics alignment is moderate globally but
weak and heterogeneous per pair: 6 of 30 pairs have negative
$\rho(C_\text{real\_z}, C_\text{real\_state})$, and the per-pair standard
deviation is 0.502.  This heterogeneity motivates the per-subset analysis in
Appendix~\ref{app:track_a} and the endpoint-planning decoupling framework in
Section~4.  These per-pair means are plotted in Figure~\ref{fig:motivation}b.

\begin{table}[h]
\centering
\caption{Three-cost attribution at PushT offset~50.  Global correlations pool
all 1{,}200 records; per-pair correlations average 30 within-pair Spearman
values ($\pm$\,SEM).}
\label{tab:three-cost}
\begin{tabular}{lcc}
\toprule
Correlation pair & Global Spearman & Per-pair Spearman (mean $\pm$ SEM) \\
\midrule
Predictor fidelity: $\rho(\Cmodel, C_\text{real\_z})$ & 0.864 & 0.779 $\pm$ 0.030 \\
Encoder-physics: $\rho(C_\text{real\_z}, C_\text{real\_state})$ & 0.669 & 0.353 $\pm$ 0.092 \\
End-to-end: $\rho(\Cmodel, C_\text{real\_state})$ & 0.596 & 0.332 $\pm$ 0.085 \\
\bottomrule
\end{tabular}
\end{table}

\subsection{Per-Pair Categorization}
\label{app:per-pair-cat}

The 30 offset-50 pairs were categorized by the number of successful actions
among their 40 candidates (Table~\ref{tab:pair-categories}).  This
categorization informed Phase~1 Track~A stratified sampling
(Appendix~\ref{app:track_a}): the 100-pair artifact was designed to cover the
full displacement$\times$rotation grid, including cells where Phase~0 found
zero-success pairs.

\begin{table}[h]
\centering
\caption{Per-pair categorization at PushT offset~50 ($n=30$ pairs, 40 actions
each).  Block displacement is the Euclidean distance between initial and goal
block positions in pixel coordinates.}
\label{tab:pair-categories}
\begin{tabular}{lccccc}
\toprule
Category & Definition & Count & Mean success & Mean displacement & Mean latent dist.\ \\
\midrule
Easy & $>$25/40 successes & 8 & 80.9\% & 2.57 & 12.50 \\
Hard & 1--25/40 successes & 15 & 21.2\% & 66.01 & 18.84 \\
Impossible & 0/40 successes & 7 & 0.0\% & 123.27 & 19.36 \\
\bottomrule
\end{tabular}
\end{table}

The ``Impossible'' pairs have zero successful actions in the 40-candidate pool,
indicating that the CEM search space does not contain solutions at these
displacements.  Their mean block displacement (123.3 pixels) is roughly 48$\times$
larger than the Easy pairs (2.6 pixels), while their mean latent distance
(19.4) is only 1.55$\times$ larger, illustrating the compression of physical
task structure in the learned latent geometry.

\subsection{Aggregate Latent Diagnostics}
\label{app:latent-diag}

Phase~0 also evaluated aggregate properties of the LeWM latent space over 100
randomly sampled trajectories (10 latent steps each, action blocks of 5;
seed~42).  Table~\ref{tab:latent-diagnostics} compares real encoder latents
with predictor-imagined latents.

\begin{table}[h]
\centering
\caption{Aggregate latent diagnostics for LeWM PushT.  Real latents use the
encoder applied to ground-truth observations; imagined latents use the
predictor rolled forward from the initial observation.  Effective rank is the
exponentiated entropy of the normalized singular-value spectrum.
Participation rank is $(\text{tr}\,\Sigma)^2 / \|\Sigma\|_F^2$.}
\label{tab:latent-diagnostics}
\begin{tabular}{lcc}
\toprule
Metric & Real latents & Imagined latents \\
\midrule
Temporal straightness (mean $\pm$ std) & 0.516 $\pm$ 0.301 & 0.531 $\pm$ 0.299 \\
Effective rank & 107.3 & 106.1 \\
Participation rank & 63.9 & 63.1 \\
Top-1 variance fraction & 3.40\% & 3.50\% \\
Top-10 variance fraction & 25.9\% & 26.2\% \\
Covariance trace & 190.0 & 190.0 \\
\bottomrule
\end{tabular}
\end{table}

Prediction error grows with rollout horizon: mean $\ell_2$ distance between
real and imagined latents rises from 1.10 at step~1 to 5.83 at step~10
(Table~\ref{tab:pred-error}).  Despite this drift, imagined latents preserve
aggregate distributional properties: effective rank, participation rank, and
covariance spectrum are nearly identical between real and imagined latents.
This confirms that the predictor does not collapse the representation, and the
planning failure is not attributable to predictor divergence in aggregate.

\begin{table}[h]
\centering
\caption{Predictor rollout error by horizon step.  $\ell_2$ distance is
computed between the real encoder latent and the predictor-imagined latent at
each step, averaged over 100 trajectories.}
\label{tab:pred-error}
\begin{tabular}{lcccccccccc}
\toprule
Step & 1 & 2 & 3 & 4 & 5 & 6 & 7 & 8 & 9 & 10 \\
\midrule
Mean $\ell_2$ & 1.10 & 1.59 & 2.17 & 2.50 & 2.87 & 3.40 & 4.00 & 4.52 & 5.14 & 5.83 \\
MSE & .008 & .018 & .048 & .058 & .072 & .101 & .136 & .168 & .215 & .271 \\
\bottomrule
\end{tabular}
\end{table}

The Phase~0 conclusion was that encoder goal geometry, not predictor drift or
planner weakness, is the primary source of long-horizon failure in PushT.  The
predictor faithfully preserves the cost landscape of the encoder, so the CEM
planner optimizes a faithful but misaligned objective.  This conclusion was
later refined by the Track~A analysis (Appendix~\ref{app:track_a}) and the
endpoint-planning decoupling framework (Section~4), which showed that the
misalignment is not a global encoder property but a local interaction between
cost geometry and CEM search dynamics.
